\section{Preliminaries}\label{sec:prelim}

Throughout this paper $\Omega\subset\R^2$ denotes a bounded open set,
$|\,\cdot\,|$ denotes Lebesgue measure on $\R^2$, $\HH^1$ denotes
one-dimensional Hausdorff measure, and $B_r(x)\subset\R^2$ is the open disc
of radius $r$ centred at $x$; we write $B_r = B_r(0)$.  The fixed
normalisation $|\Omega|=\pi$ is in force from \cref{sec:intro} onward but is
not required for the statements in this section unless explicitly noted.  We
collect here the definitions and classical results the rest of the paper relies
on.  Standard results are stated precisely with full hypotheses and cited to
published sources.

%--------------------------------------------------------------------
\subsection{Sets of finite perimeter}
%--------------------------------------------------------------------

We follow the conventions of \cite{Maggi2012}.

\begin{definition}[Perimeter]\label{def:perimeter}
Let $E\subset\R^2$ be Lebesgue measurable and let $A\subset\R^2$ be open.
The \emph{perimeter of $E$ in $A$} is
\[
  \Per(E;A) := \sup\!\left\{\,\int_E \divg\varphi\,dx \;:\;
    \varphi\in C^1_c(A;\R^2),\ \|\varphi\|_{L^\infty}\le 1\,\right\}.
\]
We set $\Per(E):=\Per(E;\R^2)$.  The set $E$ is of \emph{finite perimeter} (in
$A$) if $\Per(E;A)<\infty$.
\end{definition}

If $E$ has locally finite perimeter, the distributional gradient
$\mu_E := -D\mathbf{1}_E$ is an $\R^2$-valued Radon measure with total
variation $|\mu_E|(A)=\Per(E;A)$
\cite[Proposition~12.1, Definition~12.1]{Maggi2012}.

\begin{definition}[Reduced boundary]\label{def:reduced-boundary}
For $E$ of locally finite perimeter, the \emph{reduced boundary} $\partial^*E$
is the set of $x\in\supp|\mu_E|$ at which
\[
  \nu_E(x) := \lim_{r\to 0^+}\frac{\mu_E(B_r(x))}{|\mu_E|(B_r(x))}
\]
exists and satisfies $|\nu_E(x)|=1$.  The vector $\nu_E(x)$ is the
\emph{measure-theoretic outer unit normal} to $E$.
\end{definition}

\begin{theorem}[De~Giorgi structure theorem]\label{thm:de-giorgi}
Let $E\subset\R^2$ have locally finite perimeter.  Then $\partial^*E$ is
countably $\HH^1$-rectifiable, $\mu_E = \nu_E\,\HH^1\lfloor\partial^*E$, and
for every Borel set $A$,
\[
  \Per(E;A) = \HH^1(\partial^*E\cap A).
\]
\end{theorem}

\begin{proof}
See \cite[Theorem~15.9, Corollary~16.1]{Maggi2012} or
\cite[Section~5.7, Theorem~5.15]{EvansGariepy2015}.
\end{proof}

\begin{definition}[Essential boundary]\label{def:essential-boundary}
For a measurable set $E\subset\R^2$, its \emph{essential boundary} is
$\partial^e E := \R^2\setminus(E^{(0)}\cup E^{(1)})$, where
$E^{(t)}:=\{x:\lim_{r\to0}|E\cap B_r(x)|/|B_r(x)|=t\}$.
\end{definition}

\begin{theorem}[Federer]\label{thm:federer}
If $E\subset\R^2$ has locally finite perimeter then
$\partial^*E\subseteq E^{(1/2)}\subseteq\partial^e E$ and
$\HH^1(\partial^e E\setminus\partial^*E)=0$.  In particular,
$\Per(E;A)=\HH^1(\partial^e E\cap A)$ for every Borel $A$.
\end{theorem}

\begin{proof}
See \cite[Theorem~15.5, Theorem~16.2]{Maggi2012} or
\cite[Theorem~5.16]{EvansGariepy2015}.
\end{proof}

%--------------------------------------------------------------------
\subsection{The coarea formula and differentiation of the distribution
function}
%--------------------------------------------------------------------

\begin{theorem}[Coarea formula]\label{thm:coarea}
Let $\Omega\subset\R^2$ be open and let $u\in W^{1,1}(\Omega)$.  Then for
every nonnegative Borel function $g:\Omega\to[0,\infty]$,
\begin{equation}\label{eq:coarea}
  \int_\Omega g(x)\,|\nabla u(x)|\,dx
  = \int_{-\infty}^{\infty}\!\!\left(\int_{\{u=t\}} g\,d\HH^1\right)dt.
\end{equation}
Moreover, for Lebesgue-almost every $t\in\R$ the superlevel set $\{u>t\}$ has
finite perimeter in $\Omega$, the level set $\{u=t\}$ is $\HH^1$-rectifiable
with $\HH^1(\{u=t\})<\infty$, and
\[
  \Per(\{u>t\};\Omega) = \HH^1(\{u=t\}) \quad\text{for a.e.\ } t.
\]
\end{theorem}

\begin{proof}
The coarea formula \eqref{eq:coarea} for $W^{1,1}$ functions is
\cite[Theorem~13.1]{Maggi2012}; see also
\cite[Section~3.4, Theorem~3.13]{EvansGariepy2015}.  The perimeter identity and
rectifiability of almost every level set follow from the $BV$-coarea formula
\cite[Theorem~3.40]{AFP2000}: since
$\int_\R\Per(\{u>t\};\Omega)\,dt = \int_\Omega|\nabla u|\,dx < \infty$, the
set $\{u>t\}$ has finite perimeter for a.e.\ $t$.
\end{proof}

The following corollary, which we call the \emph{differentiation lemma} for the
distribution function, is used repeatedly in the level-set analysis.

\begin{lemma}[Differentiation of the distribution function]\label{lem:dist-deriv}
Let $\Omega\subset\R^2$ be bounded and let $u\in C^1(\Omega)\cap C(\overline\Omega)$
with $u\ge 0$.  Define $\mu(v):=|\{u>v\}|$ for $v\ge 0$.  Then $\mu$ is
non-increasing and absolutely continuous, and for Lebesgue-almost every $v>0$,
\begin{equation}\label{eq:dist-deriv}
  -\mu'(v) = \int_{\{u=v\}} \frac{1}{|\nabla u|}\,d\HH^1.
\end{equation}
In particular, the Cauchy--Schwarz inequality applied to the pair
$(1, |\nabla u|^{1/2})$ on $\{u=v\}$ gives
\begin{equation}\label{eq:cauchy-schwarz-flux}
  \Per(\{u>v\})^2 = \HH^1(\{u=v\})^2
  \le \left(\int_{\{u=v\}}|\nabla u|\,d\HH^1\right)\left(-\mu'(v)\right)
  \quad \text{for a.e.\ }v.
\end{equation}
\end{lemma}

\begin{proof}
The coarea formula \eqref{eq:coarea} with $g = |\nabla u|^{-1}\mathbf{1}_{\{\nabla u\ne 0\}}$
gives, for $0\le s < t$,
\[
  |\{s < u \le t\}\cap\{\nabla u\ne 0\}|
  = \int_s^t \left(\int_{\{u=v\}\cap\{\nabla u\ne 0\}}\frac{1}{|\nabla u|}\,d\HH^1\right)dv.
\]
Since $u\in C^1(\Omega)$, Sard's theorem (\cref{thm:sard} below) implies that
$\{u=v\}\cap\{\nabla u=0\}$ is empty for a.e.\ $v$, so the restriction to
$\{\nabla u\ne 0\}$ is superfluous for a.e.\ $v$.  Also $|\{\nabla u=0\}\cap\{s<u\le t\}|$
contributes zero to the left side (it has measure zero by a standard argument).
Hence $\mu(s)-\mu(t) = \int_s^t(-\mu')\,dv$ with $-\mu'(v) =
\int_{\{u=v\}}|\nabla u|^{-1}\,d\HH^1$ a.e.

For \eqref{eq:cauchy-schwarz-flux}: by Cauchy--Schwarz on $\{u=v\}$,
\[
  \HH^1(\{u=v\})^2
  = \left(\int_{\{u=v\}}1\,d\HH^1\right)^{\!2}
  \le \left(\int_{\{u=v\}}|\nabla u|\,d\HH^1\right)
      \left(\int_{\{u=v\}}\frac{1}{|\nabla u|}\,d\HH^1\right)
  = \left(\int_{\{u=v\}}|\nabla u|\,d\HH^1\right)(-\mu'(v)),
\]
and by \cref{thm:coarea}, $\HH^1(\{u=v\}) = \Per(\{u>v\})$ for a.e.\ $v$.
\end{proof}

\begin{remark}\label{rem:flux-coarea}
The integral $\int_{\{u=v\}}|\nabla u|\,d\HH^1$ is the \emph{flux} of the
torsion function across the level set.  When $-\Delta u = 1$ in $\{u>v\}$ with
$u=v$ on $\partial\{u>v\}$ (for a.e.\ regular $v$), integration of
$-\Delta u = 1$ over $\{u>v\}$ and Green's first identity give
$\int_{\{u=v\}}|\nabla u|\,d\HH^1 = |\{u>v\}| = \mu(v)$.  Hence
\eqref{eq:cauchy-schwarz-flux} becomes
$\Per(\{u>v\})^2\le \mu(v)\,(-\mu'(v))$ for a.e.\ $v$; the deficit from
equality is $D_H(v) := \mu(v)(-\mu'(v)) - \Per(\{u>v\})^2 \ge 0$.  This is
developed in \cref{sec:identity}.
\end{remark}

%--------------------------------------------------------------------
\subsection{Isoperimetric inequalities}
%--------------------------------------------------------------------

\begin{theorem}[Planar isoperimetric inequality]\label{thm:isoperimetric}
For every set of finite perimeter $E\subset\R^2$ with $|E|<\infty$,
\begin{equation}\label{eq:isoperimetric}
  \Per(E)^2 \ge 4\pi\,|E|,
\end{equation}
with equality if and only if $E$ is equivalent (up to a Lebesgue-null set)
to a disc $B_r(x)$.
\end{theorem}

\begin{proof}
This is \cite[Theorem~14.1]{Maggi2012}; see also
\cite[Section~5.6, Theorem~5.6]{EvansGariepy2015}.
\end{proof}

\begin{theorem}[Relative isoperimetric inequality in a disc]
\label{thm:rel-isoperimetric}
There is an absolute constant $C_{\mathrm{rel}}>0$ such that for every set of
finite perimeter $E\subset\R^2$, every $x\in\R^2$, and every $r>0$,
\begin{equation}\label{eq:rel-isoperimetric}
  \min\!\bigl\{|E\cap B_r(x)|,\,|B_r(x)\setminus E|\bigr\}
  \;\le\; C_{\mathrm{rel}}\,\Per(E;B_r(x))^2.
\end{equation}
The constant $C_{\mathrm{rel}}$ is independent of $E$, $x$, and $r$.
\end{theorem}

\begin{proof}
By the relative isoperimetric inequality for convex domains, applied to the unit
disc $B_1\subset\R^2$, there is a constant $C_{\mathrm{rel}}>0$ such that for
every set of finite perimeter $E$,
$\min\{|E\cap B_1|, |B_1\setminus E|\} \le C_{\mathrm{rel}}\,\Per(E;B_1)^2$;
see \cite[Theorem~12.37]{Maggi2012} (stated in $\R^n$, here with $n=2$,
$K=B_1$) or \cite[Theorem~5.9]{EvansGariepy2015}.  The general case at radius
$r$ follows by the substitution $\tilde{E}:=\{y/r:y\in E\}$: since
$|\tilde{E}\cap B_1| = r^{-2}|E\cap B_r(x)|$ and
$\Per(\tilde{E};B_1) = r^{-1}\Per(E;B_r(x))$, the unit-disc inequality for
$\tilde{E}$ multiplied by $r^2$ yields \eqref{eq:rel-isoperimetric} with the
same constant $C_{\mathrm{rel}}$.
\end{proof}

%--------------------------------------------------------------------
\subsection{The sharp quantitative isoperimetric inequality}
%--------------------------------------------------------------------

For a set of finite perimeter $E\subset\R^2$ with $0<|E|<\infty$ we define the
\emph{Fraenkel asymmetry}
\begin{equation}\label{eq:fraenkel}
  \AF(E) := \min_{x_0\in\R^2}\frac{|E\triangle B_{r_E}(x_0)|}{|E|},
  \qquad r_E := \sqrt{|E|/\pi},
\end{equation}
and the \emph{isoperimetric deficit}
\begin{equation}\label{eq:iso-deficit}
  \diso(E) := \frac{\Per(E)}{2\sqrt{\pi|E|}}-1 \;\ge\; 0,
\end{equation}
where $2\sqrt{\pi|E|}$ is the perimeter of a disc of area $|E|$.  Both
quantities are dimensionless, translation- and scaling-invariant, and
nonnegative, vanishing precisely when $E$ is a disc.

\begin{theorem}[Fusco--Maggi--Pratelli]\label{thm:fmp}
There is an absolute constant $C_{\mathrm{F}}>0$ such that for every set of
finite perimeter $E\subset\R^2$ with $0<|E|<\infty$,
\begin{equation}\label{eq:fmp}
  \AF(E) \;\le\; C_{\mathrm{F}}\,\diso(E)^{1/2}.
\end{equation}
The exponent $\tfrac{1}{2}$ is optimal: it cannot be replaced by any strictly
larger power, as shown by elongated ellipses with small eccentricity.
The value of $C_{\mathrm{F}}$ is not explicit in \cite{FMP2008}; it is finite
and absolute (depending on neither $E$ nor any geometric parameter of $\Omega$).
\end{theorem}

\begin{proof}
This is Theorem~1.1 of \cite{FMP2008} (stated there in $\R^n$; the planar case
$n=2$ is the specialisation $C_{\mathrm{F}}:=C(2)$).  The theorem is proved in
\cite{FMP2008} using symmetrisation and a selection-principle argument; we use
it as a black box.
\end{proof}

\begin{remark}[Relation between $\diso$ and the additive deficit]\label{rem:deficit-forms}
The additive isoperimetric deficit
$D_I(E) := \Per(E)^2 - 4\pi|E| \ge 0$ (see \cref{def:deficits}) satisfies
\[
  \diso(E) = \sqrt{1 + \frac{D_I(E)}{4\pi|E|}}-1
  = \frac{D_I(E)/\!(4\pi|E|)}{\sqrt{1+D_I(E)/(4\pi|E|)}+1}
  \;\le\; \frac{D_I(E)}{8\pi|E|},
\]
since $\sqrt{1+s}+1\ge 2$ for $s\ge 0$.
Hence $\diso(E)^{1/2} \le (D_I(E)/(8\pi|E|))^{1/2}$, and \eqref{eq:fmp} gives
\[
  \AF(E)^2 \;\le\; C_{\mathrm{F}}^2\,\diso(E)
  \;\le\; \frac{C_{\mathrm{F}}^2}{8\pi}\,\frac{D_I(E)}{|E|}.
\]
We will apply this estimate in \cref{sec:main} to the inset $\Etil$.
\end{remark}

%--------------------------------------------------------------------
\subsection{Interior regularity of the torsion function}
%--------------------------------------------------------------------

Let $u\in H^1_0(\Omega)$ be the unique weak solution of the torsion problem
\begin{equation}\label{eq:torsion}
  -\Delta u = 1 \text{ in }\Omega, \qquad u = 0 \text{ on }\partial\Omega,
\end{equation}
i.e., $\int_\Omega\nabla u\cdot\nabla\varphi\,dx=\int_\Omega\varphi\,dx$ for
all $\varphi\in H^1_0(\Omega)$.  Existence and uniqueness follow from the
Lax--Milgram theorem; $u\ge 0$ by the weak maximum principle
\cite[Theorem~8.1]{GilbargTrudinger2001}.

\begin{theorem}[Interior smoothness and analyticity of $u$]\label{thm:torsion-reg}
Let $\Omega\subset\R^2$ be a bounded open set and let $u\in H^1_0(\Omega)$ be the
weak solution of \eqref{eq:torsion}.  Then $u\in C^\infty(\Omega)$.
Moreover, $u$ is real-analytic in $\Omega$.
\end{theorem}

\begin{proof}
Since the right-hand side of \eqref{eq:torsion} is $f\equiv 1\in C^\infty(\R^2)$
and the operator $-\Delta$ has smooth (constant) coefficients, interior Schauder
estimates bootstrap to give $u\in C^{k,\alpha}_{\mathrm{loc}}(\Omega)$ for
every $k\ge 0$ and every $\alpha\in(0,1)$; see
\cite[Theorem~6.17, Corollary~6.18]{GilbargTrudinger2001}.  Hence
$u\in C^\infty(\Omega)$.

Real-analyticity follows from the general fact that solutions of linear
elliptic equations with real-analytic coefficients and real-analytic right-hand
side are real-analytic in the interior: the Laplacian has constant (hence
analytic) coefficients and $f\equiv 1$ is analytic, so $u$ is real-analytic in
$\Omega$; this is a classical result of Morrey--Nirenberg, discussed in
\cite[Chapter~6, \S6.8 and the historical notes]{GilbargTrudinger2001}.
\end{proof}

\begin{theorem}[Interior gradient and Hessian estimates]\label{thm:interior-est}
There is an absolute constant $C>0$ with the following property.  Let
$\Omega\subset\R^2$ be a bounded open set, let $u$ satisfy \eqref{eq:torsion}
weakly, let $x\in\Omega$, and set $d:=\dist(x,\partial\Omega)$.  Then
\begin{equation}\label{eq:interior-est}
  |\nabla u(x)| \;\le\; C\!\left(\frac{\|u\|_{L^\infty(B_d(x))}}{d}+d\right),
  \qquad
  |D^2 u(x)| \;\le\; C\!\left(\frac{\|u\|_{L^\infty(B_d(x))}}{d^2}+1\right).
\end{equation}
\end{theorem}

\begin{proof}
Apply the interior Schauder estimate of \cite[Theorem~6.2]{GilbargTrudinger2001}
to $u$ on $B_d(x)$ where $-\Delta u = 1$.  The scale-invariant form of the
estimate gives
\[
  d\,\|\nabla u\|_{L^\infty(B_{d/2}(x))} + d^2\,\|D^2u\|_{L^\infty(B_{d/2}(x))}
  \;\le\; C\!\left(\|u\|_{L^\infty(B_d(x))} + d^2\right),
\]
with $C$ depending only on the dimension $n=2$.  Dividing by $d$ and $d^2$
respectively yields \eqref{eq:interior-est} (the sup over $B_{d/2}(x)$ bounds
the pointwise value at $x$).
\end{proof}

\begin{lemma}[$L^\infty$ bound]\label{lem:linfty}
Let $\Omega\subset\R^2$ be bounded and let $u$ solve \eqref{eq:torsion}.
Then $0\le u(x)\le \tfrac14(\diam\Omega)^2$ for all $x\in\Omega$.
In particular $m:=\max_{\overline\Omega}u \le \tfrac14(\diam\Omega)^2$.
\end{lemma}

\begin{proof}
The lower bound $u\ge 0$ follows from the weak maximum principle
\cite[Theorem~8.1]{GilbargTrudinger2001}.  For the upper bound, let $R=\diam\Omega$
and pick any $x_0\in\overline\Omega$; then $\Omega\subset B_R(x_0)$.  The
function $w(x):=(R^2-|x-x_0|^2)/4$ satisfies $-\Delta w=1$ in $\R^2$ and
$w\ge 0$ on $\Omega$.  Hence $h:=w-u$ satisfies $-\Delta h=0$ in $\Omega$ and
$h=w\ge 0$ on $\partial\Omega$, so by the maximum principle
\cite[Corollary~3.2]{GilbargTrudinger2001}, $h\ge 0$ in $\Omega$, giving
$u\le w\le R^2/4$.
\end{proof}

%--------------------------------------------------------------------
\subsection{Sard's theorem and regular values of the torsion function}
%--------------------------------------------------------------------

\begin{theorem}[Sard]\label{thm:sard}
Let $U\subseteq\R^d$ be open and $f\in C^k(U;\R^m)$ with
$k\ge\max\{1, d-m+1\}$.  Then the set of critical values of $f$,
\[
  f\!\left(\bigl\{x\in U: \operatorname{rank}\,Df(x) < m\bigr\}\right),
\]
has Lebesgue measure zero in $\R^m$.  In particular, if $f\in C^\infty(U)$
and $m=1$, the set $f(\{\nabla f=0\})\subset\R$ is a Lebesgue-null set.
\end{theorem}

\begin{proof}
See \cite[Theorem~A.6]{Maggi2012} for the statement in the context used here;
the original source is \cite{Sard1942}.
\end{proof}

\begin{lemma}[Almost every level of $u$ is regular]\label{lem:regular-values}
Let $\Omega\subset\R^2$ be a bounded open set and let $u$ solve
\eqref{eq:torsion}.  Then for Lebesgue-almost every $v\in(0,m)$ the level $v$
is a \emph{regular value} of $u$ in $\Omega$: that is, $\nabla u(x)\ne 0$ at
every point $x\in\{u=v\}\cap\Omega$.  For every such $v$,
\begin{enumerate}[label=(\roman*)]
  \item\label{it:rv-manifold}
    the level set $\{u=v\}\cap\Omega$ is a finite, disjoint union of smooth
    closed curves;
  \item\label{it:rv-perimeter}
    the superlevel set $\{u>v\}$ is a finite-perimeter set with
    $\Per(\{u>v\}) = \HH^1(\{u=v\})$;
  \item\label{it:rv-formula}
    the differentiation formula \eqref{eq:dist-deriv} holds at $v$, with
    $1/|\nabla u|$ finite and smooth on $\{u=v\}$.
\end{enumerate}
\end{lemma}

\begin{proof}
By \cref{thm:torsion-reg}, $u\in C^\infty(\Omega)$ (indeed real-analytic).
Applying \cref{thm:sard} with $d=2$, $m=1$: the set of critical values
$u(\{\nabla u=0\}\cap\Omega)\subset\R$ has Lebesgue measure zero.  Hence for
a.e.\ $v$, every point of $\{u=v\}\cap\Omega$ is a regular point.  For such $v$:

\ref{it:rv-manifold}: The implicit function theorem gives $\{u=v\}\cap\Omega$
as a smooth $1$-dimensional embedded submanifold of $\Omega$.  Since $\Omega$
is bounded and $u\in C(\overline\Omega)$ with $u=0$ on $\partial\Omega$, the
level set $\{u=v\}$ (for $v>0$) is compactly contained in $\Omega$ and hence
consists of finitely many smooth closed curves (a smooth $1$-manifold without
boundary, compact in $\Omega$).

\ref{it:rv-perimeter}: By \cref{thm:coarea},
$\Per(\{u>v\})=\HH^1(\{u=v\})$ for a.e.\ $v$.

\ref{it:rv-formula}: For a.e.\ $v$, \eqref{eq:dist-deriv} holds and
$|\nabla u|>0$ on $\{u=v\}$, so $1/|\nabla u|$ is finite and smooth there by
\cref{thm:torsion-reg}.
\end{proof}
